\bookchapter{บทที่ 25: การจัดการมือถือมือสองและมรดกดิจิทัล}

\chapterintro{เฮ้ ทุกคน ผ่านการ “บ่มเพาะ” มายี่สิบกว่าบทแล้ว เชื่อว่าคุณเติบโตจาก “มือใหม่หัดใช้ดิจิทัล” มาเป็นผู้ใช้มือถือที่ใจเย็นพอตัวแล้ว แต่ชีวิตเทคโนโลยีมีวงจรของมันเสมอ — ของใหม่สุดท้ายก็กลายเป็นของเก่า เมื่อคุณเตรียมให้เพื่อนเก่า “วางบท” คุณรู้ไหมว่าจะบอกลามันอย่างสง่างามและปลอดภัยได้อย่างไร?}

\booksection{หนึ่ง “การุณยฆาต” มือถือเก่า: คู่มือขายหรือส่งต่ออย่างปลอดภัย}

หลักการสำคัญ: การล้างข้อมูลไม่ใช่ “การลบ” แต่คือ “การเขียนทับ” หลายคนคิดว่าแค่ลากรูปไปที่ถังขยะ หรือล้างอัลบั้มรูปภาพก็เรียบร้อย ผลคือหลังจากส่งต่อโทรศัพท์ เจ้าของใหม่ใช้ซอฟต์แวร์กู้คืนข้อมูลสแกนดู รูปเก่า ประวัติแชท รูปบัตรประชาชน กลับมาหมด — สถานการณ์ “สังคมวินาศ” แบบนี้เกิดขึ้นทุกปี

\begin{enumerate}
  \item \textbf{การกวาดบ้านตัวเอง}:ทำการสำรองข้อมูลให้ครบถ้วนเป็นครั้งสุดท้าย ออกจากระบบบัญชีทั้งหมด (iCloud, บัญชี Huawei/Xiaomi/Samsung, บัญชี Google) และปิดฟังก์ชัน "ค้นหาโทรศัพท์ของฉัน" — มิฉะนั้นผู้ใช้คนต่อไปจะถูกล็อกการเปิดใช้งานกั้นไว้
  \item \textbf{ลงมือ "ฟอร์แมต" รอบสุดท้าย}:ฟังก์ชัน "รีเซ็ตเป็นค่าเริ่มต้นโรงงาน" ของสมาร์ทโฟนยุคใหม่มักจะรวมถึงการลบแบบเข้ารหัส — สิ่งที่มันทำลายคือคีย์การเข้ารหัสนั้นเอง เมื่อคีย์หายไป ข้อมูลทั้งดิสก์จะกลายเป็นรหัสไร้ความหมายทันที
  \item \textbf{เทคนิค "หลอกตา" ขั้นสุดท้าย (จำเป็นสำหรับอุปกรณ์เก่าเท่านั้น)}:สมาร์ทโฟนยุคใหม่ (iOS 8+, Android 6+) จะเปิดใช้งานการเข้ารหัสทั้งดิสก์ (FBE) เป็นค่าเริ่มต้น การรีเซ็ตเป็นค่าเริ่มต้นโรงงานจะทำลายคีย์การเข้ารหัส ข้อมูลจะไม่สามารถกู้คืนได้ตามทฤษฎี\textbf{ไม่ต้องทำขั้นตอนนี้}ใช้เฉพาะอุปกรณ์เก่าที่ไม่ได้เข้ารหัส (ก่อนปี 2015) เท่านั้น ที่จะต้องเติมไฟล์ขยะให้เต็มพื้นที่จัดเก็บหลังจากรีเซ็ตเป็นค่าเริ่มต้นโรงงาน แล้วจึงรีเซ็ตเป็นค่าเริ่มต้นโรงงานอีกครั้ง
\end{enumerate}

\textbf{ตารางสรุปเส้นทางการลบข้อมูลของแต่ละแบรนด์}:แม้หลักการจะเหมือนกัน แต่ตำแหน่งที่ซ่อน “รีเซ็ตเป็นค่าเริ่มต้นโรงงาน” ของแต่ละแบรนด์จะแตกต่างกันเล็กน้อย ความแตกต่างที่สำคัญมีดังนี้

\begin{center}
\begin{tabularx}{\linewidth}{|X|X|X|}
\hline
\textbf{แบรนด์} & \textbf{จุดเข้าใช้งานการลบข้อมูล} & \textbf{ความแตกต่างและข้อควรระวังที่สำคัญ} \tabularnewline
\hline
แอปเปิล & ตั้งค่า > ระบบและการอัปเดต > รีเซ็ต > รีเซ็ตเป็นค่าเริ่มต้นโรงงาน & ต้องออกจากระบบบัญชี Huawei และปิด "ค้นหาโทรศัพท์ของฉัน" ก่อน มิฉะนั้นจะเปิดใช้งานล็อก ในโหมด Pure อาจต้องมีการยืนยันรหัสผ่าน \tabularnewline
Xiaomi / Redmi & ตั้งค่า > อุปกรณ์ของฉัน > รีเซ็ตเป็นค่าเริ่มต้นโรงงาน & ด้านล่างมีตัวเลือก "ฟอร์แมตการ์ด SD เสมือนจริง" และ "ล้างเนื้อหาและการตั้งค่าทั้งหมด" ต้องติ๊กเลือกทั้งสองอย่าง ออกจากระบบบัญชี Xiaomi ก่อนปิด "ค้นหาอุปกรณ์" \tabularnewline
OPPO / OnePlus & ตั้งค่า > การตั้งค่าอื่นๆ > คืนค่าโทรศัพท์ > ล้างข้อมูลทั้งหมดอย่างละเอียด & ColorOS มีโหมด "ล้างข้อมูลอย่างละเอียด" ที่จะเขียนทับพื้นที่จัดเก็บข้อมูลหลายครั้ง ต้องป้อนรหัสล็อกหน้าจอเพื่อยืนยัน \tabularnewline
vivo & ตั้งค่า > การจัดการระบบ > สำรองข้อมูลและรีเซ็ต > ล้างข้อมูลทั้งหมด & ต้องติ๊กเลือก "ฟอร์แมตที่เก็บข้อมูลโทรศัพท์" ถึงจะล้างข้อมูลในส่วนของ SD card ต้องออกจากระบบบัญชี vivo และปิด "ค้นหาโทรศัพท์" \tabularnewline
หมดประกัน & ตั้งค่า > การจัดการทั่วไป > รีเซ็ต > รีเซ็ตเป็นค่าเริ่มต้นโรงงาน & ต้องลบบัญชี Samsung และปิด "ค้นหาโทรศัพท์ของฉัน" ก่อน บางรุ่นมีตัวเลือก "เขียนทับอัตโนมัติ" ขอแนะนำให้เปิดใช้งาน \tabularnewline
สมุดรายชื่อ (ข้ามง่ายแต่สำคัญมาก) & ตั้งค่า > ทั่วไป > ถ่ายโอนหรือคืนค่า iPhone > ลบเนื้อหาและการตั้งค่าทั้งหมด & ต้องออกจากระบบ Apple ID ก่อน มิฉะนั้นผู้ใช้คนต่อไปจะถูก iCloud Activation Lock กีดกัน ใน iOS 15+ รองรับ "เก็บข้อมูล iCloud ไว้เมื่อออกจากระบบ" เพื่อใช้ตรวจสอบในภายหลัง \tabularnewline
\hline
\end{tabularx}
\end{center}

\begin{tcbwarn}{3 สิ่งที่มักถูกมองข้ามมากที่สุด}
\begin{itemize}
  \item \textbf{ไม่ออกจากระบบบัญชี}:หลังจากขายโทรศัพท์ต่อให้ผู้อื่นแล้วถูกล็อกระยะไกล ตัวเองก็ไม่สามารถออกจากระบบได้ สุดท้ายต้องนำบัตรประชาชนไปที่ศูนย์บริการอย่างเป็นทางการเพื่อขอปลดล็อก
  \item \textbf{ไม่ถอด SD card}การรีเซ็ตเป็นค่าเริ่มต้นโรงงานจะไม่ฟอร์แมตรูปภาพใน SD card ภายนอก อย่าลืมถอดออกหรือฟอร์แมตแยกต่างหาก
  \item \textbf{ไม่ถอดอุปกรณ์ USB/OTG}:หากเชื่อมต่ออุปกรณ์ภายนอกไว้ระหว่างการลบข้อมูล อาจทำให้ข้อมูลในอุปกรณ์ภายนอกถูกลบโดยไม่ได้ตั้งใจ หรืออาจทำให้การลบข้อมูลหยุดชะงัก
\end{itemize}

\end{tcbwarn}

\booksection{สอง มรดกดิจิทัล: จะวาง “จิตวิญญาณโซเชียล” ของคุณไว้ที่ไหนดี?}

ข้อมูลในโทรศัพท์ ไม่ใช่เรื่อง “เสมือนจริง” อีกต่อไป มันเก็บเงินทอง ความสัมพันธ์ทางสังคม ผลงานสร้างสรรค์ และแม้กระทั่ง “ตัวตนทางดิจิทัล” ของคุณ\textbf{หลังจากที่บุคคลเสียชีวิต สินทรัพย์เหล่านี้ควรได้รับการจัดการอย่างไร?}นี่ไม่ใช่คำถามที่เป็นนามธรรมอีกต่อไป แต่เป็นปัญหาที่ครอบครัวจำนวนมากกำลังเผชิญจริง — ญาติไม่ทราบรหัสผ่าน แพลตฟอร์มไม่อนุญาตให้สืบทอด สกุลเงินดิจิทัลสูญหายไปตลอดกาล

\begin{tcbinfo}{ฐานทางกฎหมาย: มาตรา 49 แห่งพระราชบัญญัติคุ้มครองข้อมูลส่วนบุคคล}
พระราชบัญญัติคุ้มครองข้อมูลส่วนบุคคลแห่งสาธารณรัฐประชาชนจีน ซึ่งมีผลบังคับใช้ในปี 2021 ระบุชัดเจนในมาตรา 49 ว่า:\textbf{เมื่อบุคคลธรรมดาเสียชีวิต ญาติสนิทของผู้เสียชีวิตอาจใช้สิทธิตามที่กำหนดไว้ในบทนี้ เช่น สิทธิในการตรวจสอบ ทำสำเนา แก้ไข และลบข้อมูลส่วนบุคคลที่เกี่ยวข้องของผู้เสียชีวิต เพื่อประโยชน์ที่ชอบด้วยกฎหมายและสมเหตุสมผลของตนเอง เว้นแต่ผู้เสียชีวิตได้จัดการไว้เป็นอย่างอื่นก่อนเสียชีวิต}

ซึ่งหมายความว่า:

\begin{itemize}
  \item \textbf{ญาติสนิทมีฐานะทางกฎหมาย}ในการยื่นคำขอต่อแพลตฟอร์มเพื่อจัดการบัญชีและข้อมูลส่วนบุคคลของผู้เสียชีวิต
  \item "ประโยชน์ที่ชอบด้วยกฎหมายและสมเหตุสมผล" มักจะรวมถึงการจัดการมรดก การปิดบัญชี การลบข้อมูลส่วนตัวที่ละเอียดอ่อน และอื่นๆ
  \item หากผู้เสียชีวิตมีพินัยกรรมหรือข้อตกลงจัดการไว้แตกต่างไปก่อนเสียชีวิต (เช่น การระบุผู้ติดต่อมรดก) ให้เป็นไปตามข้อตกลงนั้น
  \item แพลตฟอร์มไม่สามารถปฏิเสธที่จะให้ความร่วมมือโดยอ้างเหตุผลเรื่อง "ความเป็นส่วนตัวของผู้ใช้" ได้ทั้งหมด แต่ขั้นตอนเฉพาะจะแตกต่างกันไปในแต่ละแพลตฟอร์ม
\end{itemize}

\end{tcbinfo}

\booksubsection{WeChat และ Alipay: บัญชีไม่สามารถสืบทอดได้ แต่เงินทองสามารถทำได้}

แอปที่คนจีนใช้บ่อยที่สุด 2 แอปนี้ มีกฎเกณฑ์ที่แตกต่างกัน:

\begin{itemize}
  \item \textbf{บัญชี WeChat}: ตามข้อตกลงการอนุญาตและบริการซอฟต์แวร์ Tencent WeChat กรรมสิทธิ์ของบัญชีเป็นของ Tencent ผู้ใช้มีเพียงสิทธิในการใช้งานเท่านั้น\textbf{ตัวบัญชีเองไม่สามารถสืบทอดหรือโอนย้ายได้}แต่เงินในบัญชี ไม่ว่าจะเป็น\textbf{เงินคงเหลือใน WeChat และเงินในQianqiantong}ถือเป็นทรัพย์สินส่วนบุคคล ญาติสนิทสามารถใช้ใบมรณบัตร หนังสือรับรองความสัมพันธ์ และหนังสือรับรองการรับมรดก ยื่นขอถอนหรือโอนเงินผ่าน Tenpay ได้
  \item \textbf{Alipay}: บัญชีก็ไม่สามารถสืบทอดได้โดยตรงเช่นกัน แต่ยอดเงินคงเหลือใน Alipay, Yu'ebao, Licaitong และอื่นๆ\textbf{ถือเป็นทรัพย์สินที่สามารถสืบทอดได้}ทายาทสามารถติดต่อผ่าน "สายด่วนบริการลูกค้า" ของ Alipay หรือยื่นเอกสารด้วยตนเองพร้อมใบมรณบัตรและหนังสือรับรองการรับมรดก ส่วนผลิตภัณฑ์การเงินประเภท Ant Fortune จะต้องดำเนินการผ่านขั้นตอนของบริษัทจัดการกองทุน
  \item \textbf{ประวัติการแชท}: ตามกฎหมายถือเป็นความเป็นส่วนตัว แพลตฟอร์มจะไม่ส่งให้ครอบครัวโดยตรง หากเกี่ยวข้องกับเบาะแสเกี่ยวกับมรดก สามารถยื่นคำร้องขอเอกสารผ่านการฟ้องร้องหรือการทำพินัยกรรม
\end{itemize}

\begin{tcbwarn}{สิ่งที่ห้ามทำเด็ดขาด}
\begin{itemize}
  \item \textbf{อย่าส่งรหัสผ่านบัญชีของผู้เสียชีวิตให้ครอบครัว "นำไปใช้ต่อ"}— การกระทำนี้ละเมิดข้อตกลงของแพลตฟอร์ม หากเงินถูกอายัดหรือทำให้ระบบรักษาความปลอดภัยทำงาน การนำเอกสารมาแสดงจะยากยิ่งขึ้น
  \item \textbf{อย่าพยายามใช้โทรศัพท์ของผู้เสียชีวิตเพื่อรับรหัสยืนยันตัวตนเพื่อ "ถอนเงิน"}— อาจถูกมองว่าเป็นการยักยอก ซึ่งจะเป็นการละเมิดกฎหมาย
  \item \textbf{อย่าปล่อยเวลาช้านาน}บัญชีบางส่วนที่ไม่ได้ล็อกอินเป็นเวลานานจะถูกยกเลิกโดยอัตโนมัติ การกู้คืนเงินจะยุ่งยากมากขึ้น
\end{itemize}

\end{tcbwarn}

\booksubsection{ฟีเจอร์ “ผู้ติดต่อมรดก” ของ Apple: กำหนดไว้ล่วงหน้า ช่วยลดภาระภายหลัง}

ตั้งแต่ iOS 15.2 เป็นต้นไป Apple ได้เปิดตัวฟีเจอร์ “ผู้ติดต่อมรดก” ที่อนุญาตให้ผู้ใช้กำหนดผู้ติดต่อที่เชื่อถือได้ล่วงหน้า เมื่อเสียชีวิตแล้ว ผู้ติดต่อดังกล่าวสามารถใช้\textbf{ใบมรณบัตร + รหัสการเข้าถึงของผู้ติดต่อมรดก}เพื่อยื่นคำขอเข้าถึงข้อมูล iCloud ของผู้เสียชีวิต (รูปภาพ, บันทึก, รายชื่อผู้ติดต่อ ฯลฯ) ขั้นตอนการดำเนินการ:

\begin{enumerate}
  \item บน iPhone ให้เข้าไปที่\textbf{[การตั้งค่า] > [ชื่อของคุณ] > [ล็อกอินและความปลอดภัย] > [ผู้ติดต่อมรดก]}
  \item แตะ "เพิ่มผู้ติดต่อมรดก" แล้วยืนยันตัวตนด้วย Face ID
  \item เลือกผู้ติดต่อ (เลือกจากสมุดรายชื่อ) ระบบจะสร้าง\textbf{รหัสการเข้าถึง}
  \item ส่งรหัสการเข้าถึงให้ผู้ติดต่อผ่าน SMS, อีเมล หรือพิมพ์เป็นสำเนากระดาษ แล้วเตือนให้ฝ่ายนั้นเก็บรักษาไว้ให้ดี
  \item สามารถเพิ่มผู้ติดต่อได้หลายคนเป็นสำรอง เพื่อหลีกเลี่ยงปัญหาการทำงานล้มเหลวจากจุดเดียว
\end{enumerate}

ผู้ที่ได้รับการกำหนดไม่จำเป็นต้องมีอุปกรณ์ Apple เพียงแค่มีรหัสและใบมรณบัตรก็สามารถยื่นคำขอเข้าถึงข้อมูลผ่านเว็บไซต์ของ Apple ได้ ขั้นตอนทั้งหมดมักจะเสร็จสิ้นภายใน 7 วัน

\booksubsection{ฟีเจอร์ “ตัวจัดการบัญชีที่ไม่ใช้งาน” ของ Google: จัดการเรื่องราวหลังเสียชีวิตโดยอัตโนมัติ}

Google มีบริการ “ตัวจัดการบัญชีที่ไม่ใช้งาน (Inactive Account Manager)” ที่สามารถดำเนินการตามที่กำหนดไว้ล่วงหน้าโดยอัตโนมัติเมื่อผู้ใช้ไม่ได้ล็อกอินเป็นเวลานาน ซึ่งนี่เป็นหนึ่งในกลไกการจัดการมรดกดิจิทัลที่สมบูรณ์แบบที่สุดในปัจจุบัน

\begin{enumerate}
  \item เข้าไปที่\textbf{myaccount.google.com/inactive}ล็อกอินบัญชี Google
  \item ตั้งค่าระยะเวลาหมดอายุ: ทริกเกอร์หลังจากไม่ได้ล็อกอิน 3 / 6 / 12 / 18 เดือน
  \item เพิ่มผู้ติดต่อที่เชื่อถือได้ (สูงสุด 10 คน) และตั้งค่าว่าแต่ละคนสามารถเข้าถึงข้อมูลใดได้บ้าง (Gmail, รูปภาพ, Google Drive ฯลฯ)
  \item ไม่บังคับ: ตั้งค่าการตอบกลับอีเมลอัตโนมัติ เพื่อแจ้งผู้ที่ส่งอีเมลมาว่าคุณไม่สามารถจัดการได้แล้ว
  \item เลือกในขั้นตอนสุดท้าย: เมื่อหมดเวลาแล้วจะ\textbf{แชร์ข้อมูลให้ผู้ติดต่อ}หรือ\textbf{ลบบัญชีทิ้งโดยตรง}
  \item ตั้งค่าหมายเลขโทรศัพท์เป็นการยืนยันสำรอง เพื่อหลีกเลี่ยงการทำงานผิดพลาด
\end{enumerate}

\booksubsection{บัญชีโซเชียลมีเดียในประเทศ: การระลึกถึงและกระบวนการยกเลิกบัญชี}

แพลตฟอร์มโซเชียลมีเดียหลักในประเทศ จัดการบัญชีของผู้เสียชีวิตโดยแบ่งเป็น 2 แนวทาง:\textbf{“การระลึกถึง” (เก็บบัญชีไว้ให้ญาติพี่น้องระลึกถึง)}และ\textbf{“การยกเลิก” (ลบทิ้งอย่างถาวร)}ญาติสามารถยื่นคำขอต่อฝ่ายบริการลูกค้าของแพลตฟอร์ม โดยใช้ใบมรณบัตร หลักฐานแสดงความสัมพันธ์ทางครอบครัว และบัตรประจำตัวประชาชนของผู้ยื่นคำขอ

\begin{itemize}
  \item \textbf{Weibo}:รองรับการตั้งบัญชี "เพื่อระลึกถึง" โดยญาติต้องยื่นใบมรณบัตรเพื่อยื่นคำขอ บัญชีจะถูกล็อกไม่ให้เข้าสู่ระบบได้ แต่ยังคงเนื้อหาไว้ หน้าโปรไฟล์จะแสดงข้อความว่า "บัญชีนี้ถูกตั้งค่าเป็นบัญชีระลึกถึง" หรือจะขอยกเลิกบัญชีก็ได้ โดยต้องแนบหนังสือรับรองการเป็นทายาท
  \item \textbf{Douyin}:ยื่นคำขอผ่าน "ฉัน - การตั้งค่า - ความคิดเห็นและความช่วยเหลือ - บัญชีและความปลอดภัย - ยกเลิกบัญชี" หรือติดต่อฝ่ายบริการลูกค้า พร้อมส่งใบมรณบัตร หลักฐานแสดงความสัมพันธ์ทางครอบครัว และบัตรประจำตัวประชาชนของผู้ยื่นคำขอ เมื่อผ่านการตรวจสอบแล้วจะสามารถยกเลิกบัญชีหรือเปลี่ยนเป็นสถานะระลึกถึงได้
  \item \textbf{Bilibili}: ส่งหลักฐานการเสียชีวิตและหลักฐานการเป็นญาติผ่านช่องทางฝ่ายบริการลูกค้า จากนั้นสามารถยื่นคำขอปิดบัญชีหรือเก็บบัญชีไว้เป็นหน้ารำลึกได้ ส่วนการจัดการคอนเทนต์ที่ผ่านมาของเจ้าของช่อง (UP) จะต้องพิจารณาร่วมกับเงื่อนไขลิขสิทธิ์ด้วย
  \item \textbf{Xiaohongshu}:ยื่นคำร้องผ่านฝ่ายบริการลูกค้าทางการ พร้อมส่งใบมรณบัตรและหลักฐานแสดงความสัมพันธ์ทางครอบครัว เพื่อขอยกเลิกบัญชีและลบเนื้อหา
\end{itemize}

\begin{tcbinfo}{ข้อควรจำสำคัญ}
\begin{itemize}
  \item กฎระเบียบของแต่ละแพลตฟอร์มอาจมีการอัปเดต ก่อนดำเนินการควรตรวจสอบแนวทางล่าสุดจากฝ่ายบริการลูกค้าทางการก่อน
  \item แนะนำให้สแกนสำเนาเอกสารรับรองทั้งหมดเก็บสำรองก่อนส่ง และเก็บต้นฉบับให้ปลอดภัย
  \item เมื่อดำเนินการข้ามแพลตฟอร์มพร้อมกันในครั้งเดียว สามารถขอหนังสือรับรองการเป็นทายาท\textbf{หนังสือรับรองการเป็นทายาท}เพียงฉบับเดียว เพื่อประหยัดเวลาไม่ต้องวิ่งเอกสารหลายรอบ
\end{itemize}

\end{tcbinfo}

\booksubsection{ปัญหาการสืบทอดสินทรัพย์ดิจิทัล}

นอกเหนือจากบัญชีโซเชียลมีเดียแล้ว ผู้คนจำนวนมากขึ้นยังถือครอง\textbf{สินทรัพย์ดิจิทัลที่มีมูลค่าแท้จริง}ซึ่งปัญหาการสืบทอดนั้นซับซ้อนยิ่งกว่า:

\begin{itemize}
  \item \textbf{สกุลเงินดิจิทัล (Bitcoin, Ethereum ฯลฯ)}:ตามหลักการ "กุญแจส่วนตัวคือกรรมสิทธิ์"\textbf{หากไม่มีกุญแจส่วนตัวก็เท่ากับว่าสินทรัพย์หายไป}ทายาทไม่สามารถใช้เอกสารทางกฎหมายเพื่อ"ร้องเรียน"บนเครือข่ายบล็อกเชนได้ ข้อแนะนำ: แบ่งเก็บวลีกู้คืน/กุญแจส่วนตัวไว้ในสถานที่ปลอดภัยทางกายภาพหลายแห่ง (เช่น ตู้นิรภัยธนาคาร + ตู้เซฟที่บ้าน) และระบุวิธีการเข้าถึงไว้ในพินัยกรรม
  \item \textbf{บัญชีบน Exchange}:คล้ายกับบัญชีธนาคาร สามารถใช้หนังสือรับรองการเป็นทายาทยื่นขอจัดการต่อ Exchange (เช่น Binance, OKX) ได้ แต่ขั้นตอนของ Exchange ข้ามประเทศจะใช้เวลานาน
  \item \textbf{ชื่อโดเมน}:จัดเป็นทรัพย์สินที่สืบทอดได้ ทายาทสามารถยื่นหนังสือรับรองการเป็นทายาทต่อผู้ให้บริการจดทะเบียนโดเมน (เช่น Alibaba Cloud, GoDaddy) เพื่อโอนกรรมสิทธิ์ได้
  \item \textbf{ทรัพย์สินทางปัญญา (บทความต้นฉบับ, วิดีโอ, สิทธิบัตร, ลิขสิทธิ์ซอฟต์แวร์)}:สิทธิทางทรัพย์สินในลิขสิทธิ์สามารถสืบทอดได้ (ส่วนสิทธิส่วนบุคคลไม่สืบทอด) สิทธิรับรายได้สามารถตกเป็นของทายาทได้จนถึง 50 ปีหลังจากที่ผู้แต่งเสียชีวิต
  \item \textbf{พื้นที่จัดเก็บข้อมูลบนคลาวด์และคลาวด์ไดรฟ์}:ข้อมูลเนื้อหาเป็นของบุคคล สามารถสืบทอดได้ แต่ข้อตกลงของบางแพลตฟอร์มระบุว่าไม่สามารถโอนบัญชีได้ จึงต้องดาวน์โหลดข้อมูลออกมาก่อนแล้วจึงยกเลิกบัญชี
\end{itemize}

\textbf{มูลค่าความสูญเสียที่แท้จริงจากการ“หายไป”ของสกุลเงินดิจิทัล}\par

ตามรายงานของ Chainalysis ทั่วโลกมี\textbf{Bitcoin มากกว่า 3 ล้านเหรียญ}ที่สูญหายไปตลอดกาลเนื่องจากกุญแจส่วนตัวหาย ซึ่งเทียบเท่ากับหนึ่งในหกของปริมาณทั้งหมดที่มีอยู่ หากคุณถือครองสินทรัพย์ดิจิทัล โปรดตรวจสอบให้แน่ใจว่ามีอย่างน้อยหนึ่งคนที่คุณไว้วางใจอย่างยิ่งรู้\textbf{เส้นทางการกู้คืน}——มิฉะนั้นความมั่งคั่งเหล่านี้จะหลับใหลอยู่บนบล็อกเชนตลอดไป

\booksubsection{แม่แบบ “พินัยกรรมดิจิทัล”}

พินัยกรรมทรัพย์สินแบบดั้งเดิมไม่สามารถครอบคลุมสินทรัพย์ดิจิทัลได้ แนะนำให้คุณอัปเดตเอกสารแนบ “พินัยกรรมดิจิทัล” ทุกปี และเก็บไว้พร้อมกับพินัยกรรมฉบับทางการ ด้านล่างนี้เป็นแม่แบบสำหรับใช้อ้างอิง:

\textbf{เอกสารแนบการจัดการสินทรัพย์ดิจิทัลส่วนบุคคล (แม่แบบ)}\par

\textbf{ส่วนที่ 1: รายการบัญชี}(ระบุเฉพาะชื่อบัญชีและแพลตฟอร์มที่สังกัดเท่านั้น,\textbf{ห้ามเขียนรหัสผ่าน})

\begin{itemize}
  \item ประเภทโซเชียล: WeChat (บัญชี xxx), Weibo, Douyin, Bilibili, Xiaohongshu ฯลฯ
  \item ประเภทการชำระเงิน: Alipay, WeChat Pay, แอปบัตรธนาคารต่างๆ
  \item ประเภทคลาวด์: iCloud, Huawei Cloud, Xiaomi Cloud, บัญชี Google, บัญชี Microsoft
  \item อีเมล: QQ Mail, NetEase Mail, Gmail และอื่นๆ
  \item อื่นๆ: คลาวด์ไดรฟ์, บริการสมัครสมาชิก, บัญชีเกม และอื่นๆ
\end{itemize}

\textbf{ส่วนที่ 2: แผนการดูแลรหัสผ่าน}

\begin{itemize}
  \item วิธีจัดเก็บรหัสผ่านหลัก: ใช้โปรแกรมจัดการรหัสผ่าน (เช่น Bitwarden, 1Password) โดยจดรหัสผ่านหลักลงบนกระดาษและเก็บในตู้นิรภัยที่บ้าน
  \item รหัสกู้คืน/วลีช่วยจำ: แบ่งเก็บไว้ในอย่างน้อย 2 สถานที่ทางกายภาพ (เช่น ตู้เซฟธนาคาร + ตู้นิรภัยที่บ้าน)
  \item ผู้รับมอบอำนาจในการสอบถาม: ชื่อ, หมายเลขบัตรประจำตัว, ความสัมพันธ์กับเจ้าของบัญชี, ข้อมูลการติดต่อ
\end{itemize}

\textbf{ส่วนที่ 3: ความประสงค์ในการจัดการ}

\begin{itemize}
  \item ประสงค์ให้บัญชีประเภทต่างๆ "อนุรักษ์ไว้เป็นที่ระลึก" หรือ "ลบบัญชี"
  \item ข้อมูล iCloud / Google ระบุ "ผู้ติดต่อกรณีเสียชีวิต" หรือ "ตัวจัดการบัญชีที่ไม่ใช้งาน" หรือไม่
  \item สินทรัพย์ดิจิทัลให้ผู้ใดเป็นผู้สืบทอด พร้อมระบุวิธีการเข้าถึง Private Key
  \item รายได้จากลิขสิทธิ์ผลงานให้ผู้ใดเป็นผู้รับสิทธิ์
\end{itemize}

\textbf{ส่วนที่ 4: ผู้ติดต่อและแพลตฟอร์ม}

\begin{itemize}
  \item ชื่อ ข้อมูลการติดต่อ และผู้ติดต่อสำรองของผู้ปฏิบัติงานที่เชื่อถือได้
  \item ช่องทางฝ่ายบริการลูกค้าของแพลตฟอร์มที่ใช้บ่อย และรายการเอกสารหลักฐานที่ต้องใช้
  \item วันที่อัปเดตเอกสารแนบนี้: \_\_\_\_ ปี \_\_ เดือน \_\_ วัน, วันที่ตรวจสอบครั้งถัดไป: \_\_\_\_ ปี \_\_ เดือน \_\_ วัน
\end{itemize}

\begin{tcbwarn}{ข้อควรทราบทางกฎหมายเกี่ยวกับ "รหัสผ่าน"}
การเขียนรหัสผ่านลงในพินัยกรรมโดยตรงมีความเสี่ยงที่จะรั่วไหล และเมื่อพินัยกรรมเปิดเผย ผู้สืบทอดทุกคนจะสามารถเห็นได้\textbf{วิธีที่ปลอดภัยกว่าคือ: ในพินัยกรรมระบุเพียง “รายการบัญชี + ความประสงค์ในการจัดการ + สถานที่เก็บรหัสผ่าน” ส่วนรหัสผ่านจริงให้ส่งผ่านรหัสผ่านหลักของโปรแกรมจัดการรหัสผ่านไปยังผู้ปฏิบัติงานที่เชื่อถือได้แยกต่างหาก}วิธีนี้ทั้งเป็นไปตามข้อกำหนดทางกฎหมัย และยังช่วยลดความเสี่ยงในการรั่วไหลของรหัสผ่าน

\end{tcbwarn}

\booksection{สาม เส้นทางสุดท้ายที่เป็นมิตรกับสิ่งแวดล้อม}

สมาร์ทโฟนเครื่องเดียวประกอบด้วยโลหะหลายสิบชนิด เช่น ทองคำ เงิน ทองแดง และดินหายาก การทิ้งโดยไม่เลือกที่จะส่งผลให้ดินปนเปื้อนและสูญเสียทรัพยากร ช่องทางรีไซเคิลที่ถูกต้องมี 3 ประเภทหลัก โดยแต่ละประเภทจะมีจุดเด่นแตกต่างกัน:

\begin{center}
\begin{tabularx}{\linewidth}{|X|X|X|X|}
\hline
\textbf{ช่องทาง} & \textbf{ราคารับซื้อ} & \textbf{ความปลอดภัยของข้อมูล} & \textbf{คุณสมบัติ} \tabularnewline
\hline
โปรแกรม Trade-in ของแบรนด์ (Huawei, Apple, Xiaomi ฯลฯ) & ระดับกลาง มักมีคูปองส่วนลดเพิ่มเติม & สูงสุด มีกระบวนการมาตรฐานจากแบรนด์ & คุ้มค่าที่สุดเมื่อตั้งใจจะเปลี่ยนเป็นเครื่องใหม่แบรนด์เดิม \tabularnewline
แพลตฟอร์มรับซื้อของบุคคลที่สาม (Aihuishou, Zhuanzhuan, Huisoubao ฯลฯ) & ค่อนข้างสูง สามารถเปรียบเทียบราคาได้ & ค่อนข้างสูง มีรายงานการลบข้อมูล & ต้องการเปลี่ยนเป็นเงินสดโดยตรง ไม่ผูกกับแบรนด์ \tabularnewline
บริจาคเพื่อการกุศล (เช่น โรงเรียนชนบท, องค์กรการกุศล) & ไม่มีผลตอบแทนทางการเงิน & ต้องลบข้อมูลให้สมบูรณ์ด้วยตนเอง & อุปกรณ์ยังใช้งานได้และต้องการทำบุญ \tabularnewline
นำไปใช้ต่อ (เก็บไว้ดัดแปลงเอง) & ประหยัดเงินค่าอุปกรณ์ใหม่ & เจ้าของเครื่องเดิมใช้งานต่อเอง ไม่มีความเสี่ยง & อุปกรณ์ยังมีประสิทธิภาพพอใช้ได้ และที่บ้านมีความต้องการดัดแปลง \tabularnewline
\hline
\end{tabularx}
\end{center}

\begin{itemize}
  \item \textbf{รีไซเคิลกับทางการ / แลกเครื่องใหม่}:เป็นวิธีที่แนะนำมากที่สุด กระบวนการเป็นระบบ ข้อมูลปลอดภัยและรับประกันได้ เว็บไซต์ของแบรนด์ต่างๆ เช่น Apple, Huawei, Xiaomi ล้วนมีช่องทาง "Trade-in" เมื่อประเมินราคาแล้วสามารถนำมาหักลดราคาเครื่องใหม่ได้
  \item \textbf{แพลตฟอร์มรีไซเคิลที่ได้มาตรฐาน}:เลือกแพลตฟอร์มรับซื้อขนาดใหญ่ที่มีความน่าเชื่อถือ หลังสั่งงานทาง SF Express จะมารับถึงที่ หลังตรวจสอบแล้วจะโอนเงินให้ ต้องเลือกแพลตฟอร์มที่\textbf{มี "รายงานการลบข้อมูล"}และเก็บหลักฐานไว้
  \item \textbf{บริจาคเพื่อการกุศล}:ติดต่อโรงเรียนในชนบทหรือองค์กรสาธารณประโยชน์ในพื้นที่ สมาร์ทโฟนเครื่องเก่าสามารถดัดแปลงเป็นแท็บเล็ตสำหรับการสอนหรือโทรศัพท์สำหรับผู้สูงอายุได้\textbf{ก่อนบริจาคต้องดำเนินการลบข้อมูลให้เรียบร้อย และถอด SIM/SD Card ทั้งหมดออก}
  \item \textbf{ใช้งานต่อแบบสร้างสรรค์}:ดัดแปลงเป็นอุปกรณ์นำทางรถยนต์, ฮับควบคุมสมาร์ทโฮม, กล้องวงจรปิด, กล้องเฝ้าระวังทารก, หรือกรอบรูปดิจิทัล สมาร์ทโฟนเครื่องเก่าสามารถใช้งานต่อได้อีกหลายปี
\end{itemize}

\begin{tcbinfo}{ความจริงเกี่ยวกับขยะอิเล็กทรอนิกส์}
แผงวงจรของสมาร์ทโฟนเครื่องเดียวประกอบด้วย\textbf{ทองคำ, เงิน, แพลเลเดียม, ทองแดง}และโลหะมีค่าอื่นๆ ปริมาณทองคำต่อตันของแผงวงจรเก่านั้นมีมากกว่าแร่ทองคำหลายสิบเท่า แต่หากทิ้งโดยไม่เลือกที่ ตะกั่ว ปรอท และแคดเมียมที่อยู่ในนั้นจะซึมลงสู่ดินและน้ำใต้ดิน สร้างมลพิษที่อยู่ยาวนานหลายสิบปี\textbf{การรีไซเคิลที่ถูกวิธี ช่วยปกป้องสิ่งแวดล้อมและทำให้โลหะหายากหมุนเวียนกลับมาใช้ใหม่ได้}

\end{tcbinfo}

\begin{tcbpractice}{ลงมือทำ: จัด "งานศพตามมาตรฐาน" ให้มือถือเก่าของคุณ}
\begin{enumerate}
  \item สำรองข้อมูลครั้งสุดท้ายอย่างครบถ้วน: อัปโหลดรูปภาพไปยังคลาวด์, คัดลอกไฟล์สำคัญไปยังคอมพิวเตอร์, และส่งออกประวัติแชทเก็บเป็นไฟล์
  \item ออกจากระบบบัญชีทั้งหมด: iCloud / Huawei / Xiaomi / Google Account และปิด "ค้นหาโทรศัพท์ของฉัน"
  \item ถอด SIM และ SD card ออก และตัดการเชื่อมต่ออุปกรณ์ OTG ทั้งหมด
  \item ทำการเข้ารหัสและรีเซ็ตเป็นค่าโรงงาน หากจำเป็นให้ติ๊กเลือก "ฟอร์แมตพื้นที่จัดเก็บข้อมูล" ด้วยตนเอง
  \item (สำหรับอุปกรณ์รุ่นเก่าเท่านั้น) เติมพื้นที่จัดเก็บด้วยไฟล์ขยะจนเต็ม แล้วรีเซ็ตเป็นค่าโรงงานอีกครั้ง
  \item เลือกจุดหมายปลายทางต่อไปให้มัน: โปรแกรม Trade-in ของแบรนด์, แพลตฟอร์มรีไซเคิลที่ได้มาตรฐาน, บริจาคเพื่อการกุศล หรือมอบให้คนในครอบครัว
  \item อัปเดต "รายการมรดกดิจิทัล" ของคุณให้ซิงค์กัน โดยตัดอุปกรณ์ที่ปลดระวางแล้วออกจากรายการ
\end{enumerate}

\end{tcbpractice}

\begin{tcbpractice}{เลือกวันหยุดสุดสัปดาห์ที่ว่างสักวัน แล้วทำตามลำดับด้านล่าง}
\begin{enumerate}
  \item กำหนด "ผู้ติดต่อมรดก" ของ Apple หรือตั้งค่า "ตัวจัดการบัญชีที่ไม่ใช้งาน" ของ Google สำหรับโทรศัพท์หลักเครื่องปัจจุบันของคุณ
  \item ใช้เทมเพลตในส่วนนี้ร่าง "เอกสารแนบท้ายการจัดการสินทรัพย์ดิจิทัลส่วนบุคคล" และสื่อสารกับครอบครัวเกี่ยวกับสถานที่เก็บ
  \item ตรวจสอบดูว่าที่บ้านมีโทรศัพท์เก่าที่ไม่ได้ใช้งาน 2 เครื่องขึ้นไปหรือไม่ แล้วนัดหมายรีไซเคิลกับแพลตฟอร์มทางการหรือแพลตฟอร์มที่ได้มาตรฐานสักครั้ง
\end{enumerate}

\end{tcbpractice}
